\documentclass[11pt,a4paper]{article}
\usepackage{amsmath,amssymb,amsthm}
\usepackage[margin=2.5cm]{geometry}
\usepackage{hyperref}
\title{Structure Polynomial}
\date{}
\begin{document}
\maketitle
\section{Definition}
The \textbf{Structure Polynomial} is a polynomial that characterizes the structure of a variable elliptic system. It is defined in the context of elliptic partial differential equations, where the spectral parameter $\tau$ is expressed in terms of real and imaginary components. Specifically, the structure polynomial is associated with the complex root of the spectral parameter, denoted as:

\begin{equation}
\tau = p + iq,
\end{equation}

where $q > 0$ and $p = \text{Re}(\tau)$, $q = \text{Im}(\tau)$. The structure polynomial encapsulates the coefficients of the elliptic system, typically denoted as $\alpha$ and $\beta$, which are crucial for determining the ellipticity of the system.

\section{Examples}
Consider a variable elliptic system represented by the equations:

\begin{equation}
\begin{aligned}
u_x - \alpha v_y &= 0, \\
v_x + u_y - \beta v_y &= 0,
\end{aligned}
\end{equation}

where $\alpha(x, y)$ and $\beta(x, y)$ are functions of the spatial variables. The principal symbol of this system can be expressed as:

\begin{equation}
P(\xi, \eta) = \begin{pmatrix}
\xi & -\alpha \eta \\
\eta & \xi - \beta \eta
\end{pmatrix}.
\end{equation}

The determinant of the principal symbol, which is essential for establishing ellipticity, is given by:

\begin{equation}
\text{det}(P(\xi, \eta)) = \xi^2 - \beta \xi \eta + \alpha \eta^2.
\end{equation}

The structure polynomial arises from this determinant, and its properties dictate the behavior of the elliptic system.

\section{Key Results}
A significant result concerning the structure polynomial is its role in determining the rigidity and flatness of the elliptic structure. Specifically, the following definitions are crucial:

\begin{definition}[Rigid Structure]
The elliptic structure is called rigid if the first-order quantity $T$ defined as:

\begin{equation}
T := \tau_x + \tau \tau_y \equiv 0 \text{ on } \Omega
\end{equation}

holds true, where $\Omega$ is the domain of interest. This condition implies that the spectral parameter $\tau$ is constant along characteristic directions.

\end{definition}

\begin{definition}[Flat Structure]
The elliptic structure is termed flat if the Gaussian curvature $K$ associated with the canonical metric is zero:

\begin{equation}
K \equiv 0 \text{ on } \Omega.
\end{equation}

Both conditions depend solely on the structure coefficients $\alpha$ and $\beta$ and their derivatives, as established by Alayón-Solarz (2026) \cite{2601.19274}.

\section{Diagrams / Visual Description}
The structure polynomial can be visualized in the context of the principal symbol's determinant. The behavior of the determinant in the $(\xi, \eta)$ plane can be depicted as a surface, where the regions of ellipticity correspond to areas where the determinant is non-zero. This surface can be analyzed to understand the transitions between elliptic and non-elliptic regions.

\section{Connections}
The concept of the structure polynomial is closely related to several areas in the study of elliptic systems, including:

- **Elliptic Regularity**: The structure polynomial plays a critical role in establishing regularity results for solutions to elliptic equations.
- **Transport Equations**: The evolution of the structure coefficients is governed by transport equations, which are influenced by the structure polynomial.
- **Complex Analysis**: The relationship between the structure polynomial and complex variables is fundamental in understanding the behavior of solutions to elliptic systems.

These connections highlight the importance of the structure polynomial in both theoretical and applied mathematics, particularly in the analysis of variable elliptic systems.

\section{References}
\begin{itemize}
    \item Alayón-Solarz, D. (2026). Variable Elliptic Structures on the Plane: Transport Dynamics, Rigidity, and Function Theory. arXiv:2601.19274.
\end{itemize}\end{document}
